\cotefile{79}
\scanpage{49}

\unsure{Â|À} vrai dire, pour définir les structures triangulées $\Sigma(i)$ $(i\in\unsure{\Pi|\pi})$, il suffirait de connaître (118) sur l'ens

\[
\underset{\substack{\cap\\ S\times S}}{\widetilde A}
=\{(\deleted{\unsure{x}}s,\deleted{\unsure{y}}t)\in S\times S\mid \deleted{\unsure{x}}s\mathbin{\wedge}\deleted{\unsure{y}}t\in\Pi\}
\tag*{(124)}
\]

en tant qu'application

\[
\widetilde A\longrightarrow\Pi\qquad\deleted{\text{\unsure{applicat|appliquée}}}
\tag*{(125)}
\]

d'une partie \inserted{$\widetilde A$} de $S\times S$ dans $\Pi$, \deleted{l'image} \inserted{telle} \deleted{\unsure{inverse}} \inserted{partie} ne rencontre pas la diagonale et est \deleted{stable} stable par la symétrie $(s,t)\mapsto(t,s)$, donc sa donnée équivaut à celle d'une partie $A$ de $\unsure{\mathscr P|\mathcal P|\mathfrak P}_2(S)$, et la donnée de (125), \inserted{application symétrique,} à celle d'une application

\[
 A\xrightarrow{t}\Pi\qquad t=\text{``type''}
 \tag*{(126)}
\]

dont les fibres sont donc les $A_{\Sigma(i)}$, et on \unsure{aura|a}

\[
 A=\bigcup_{i\in\Pi}\bigl(A_{\Sigma(i)}={\unsure{\wedge|t}}^{-1}(\{i\})\bigr)
 \tag*{(127)}
\]

On peut dire que la donnée d'une famille, indexée par $\Pi$, de structures triangulées sur $S$, déterminées par les structures de graphes correspondantes via (121), équivaut à celle d'une partie $A$ de $\mathscr P_2(S)$ (\unsure{ie.|i.e.} d'une partie $\widetilde A$ de $S\times S$ ne rencontrant pas la diag. et stable par symétrie) et d'une application (surjective) (127), de façon que les fibres $A_i=A_{\Sigma(i)}$

% archive sheet number: 48
